% concatenated from main.tex
\documentclass[11]{article}
\usepackage{amsmath, amssymb, graphicx, amsthm, mathrsfs}
\usepackage{mathptmx}
\usepackage[italic]{mathastext}
\usepackage{mathtools} % For better math formatting
\usepackage[utf8]{inputenc}
\usepackage{enumitem}  % For customizing itemize environment
\usepackage[colorlinks=true]{hyperref}
\usepackage{geometry}
\geometry{margin=1in}


%\newtheorem{stat}{Statement}[section]
\newtheorem{definition}{Definition}[section]
\newtheorem{theorem}[definition]{Theorem}
\newtheorem{lemma}[definition]{Lemma}
\newtheorem{remark}[definition]{Remark}
\newtheorem{example}[definition]{Example}
\newtheorem{proposition}[definition]{Proposition}
\newtheorem{corollary}[definition]{Corollary}
\newtheorem{assumption}[definition]{Assumption}
\numberwithin{equation}{section}
%
\renewcommand{\proofname}{\textsc{Proof}}


\title{Common Fixed Points for Weakly Compatible Mappings via Bivariate Auxiliary Functions}
\author{Babu G.V.R\textsuperscript{1}, Alemayehu G. Negash\textsuperscript{2,*}, Meaza F. Bogale \textsuperscript{3}\\
\small \textsuperscript{1}Department of Mathematics, Andhra University, India\\
%\small \textsuperscript{*}Visakhapatnam-530 003, Andhra Pradesh, India. \\
\small \textsuperscript{1}Email: \texttt{gvr\_babu@hotmail.com} \\
\small \textsuperscript{2,3}Department of Mathematics, Hampton University, USA\\
$^*$ \textit{Corresponding author.} \\
%\small \textsuperscript{\dag}Visakhapatnam-530 003, Andhra Pradesh, India. \\
\small \textsuperscript{\dag}Email: \texttt{alemayehu.negash@hamptonu.edu}\\
\small \textsuperscript{\ddag}Email: \texttt{meaza.bogale@hamptonu.edu}
}
\bibliographystyle{plain}

\begin{document}

\maketitle

\begin{abstract}
We establish common fixed point theorems for two pairs of weakly compatible self-mappings using an auxiliary function of two variables. Unlike classical results, our theorems do not assume continuity of the mappings and require completeness of only one image set. The use of two-variable auxiliary functions allows us to unify and extend various existing fixed point theorems in metric spaces.


\textbf{Key words and phrases:} Weakly compatible maps, auxiliary function, common fixed point.

\textbf{AMS(2000) Mathematics Subject Classification:} 47H10, 54H25.
\end{abstract}

\section{Introduction}
The method of altering distances between points, introduced by Khan et al. \cite{Khan1984}, initiated a new approach in fixed point theory. This idea was later generalized by numerous researchers [\cite{Babu2004,Babu2007,Babu2001,Babu2007b}, \cite{Choudhury2005,Choudhury2000}, \cite{Naidu2001,Naidu2003}, \cite{Park1980}, \cite{Pathak1994}, \cite{Sastry1999,Sastry1998,Sastry2003,Sastry2000}]. In particular, Choudhury and Dutta  \cite{Choudhury2000} introduced altering distance functions involving two variables and obtained important fixed point results in metric spaces. These generalized functions serve as auxiliary tools in establishing contraction-type conditions.


\begin{definition}[\cite{Khan1984}]
A function $\phi:\mathbb{R}_{+}\to \mathbb{R}_{+}$ is called an altering distance function if the following properties are satisfied:
\begin{enumerate}
    \item[(a)] $\phi(t)=0$ if and only if $t=0$; and
    \item[(b)] $\phi$ is continuous and monotonically increasing; where $\mathbb{R}_{+}=[0,\infty)$.
\end{enumerate}
\end{definition}



\begin{definition}[\cite{Choudhury2000}]
A function $\psi:\mathbb{R}_{+}\times \mathbb{R}_{+}\to \mathbb{R}_{+}$ is said to satisfy Condition-A if
\begin{enumerate}
    \item[(i)] $\psi$ is continuous;
\end{enumerate}
\end{definition}



%\section*{Auxiliary Functions of Two Variables}

\begin{enumerate}
    \item[(ii)] $\psi$ \textit{is monotone increasing in each of its arguments;}

    \item[(iii)] $\psi(0,0)=0$ \textit{and} $\psi(t,0)=0$ \textit{implies} $t=0$.
\end{enumerate}

Throughout this paper we denote

\[
\Psi=\{\psi\,|\,\psi:\mathbb{R}_{+}\times\mathbb{R}_{+}\to\mathbb{R}_{+} \text{ satisfies Condition-A}\}
\]

We call an element $\psi\in\Psi$ \textit{an auxiliary function of two variables}.

\begin{remark} \label{rem:1.3}
If $\psi(s,t)=0$, then $s=0$.
\end{remark}

\begin{example}
The following are examples of $\psi$.

Let $\psi:\mathbb{R}_{+}\times\mathbb{R}_{+}\to\mathbb{R}_{+}$ defined by

\begin{enumerate}
    \item[(i)] $\psi(s,t)=s^{p}+t^{q}$, where $p>0$, $q>0$,
    
    \item[(ii)] $\psi(s,t)=s^{p}t^{q}+t^{r}$, where $p>0$, $q\geq 0$, $r>0$,
    
    \item[(iii)] $\psi(s,t)=\max\{s^{p},t^{q}\}$, where $p>0$, $q>0$, and
    
    \item[(iv)] $\psi(s,t)=(s^{p}+r\cdot t^{q})^{\lambda}$, where $p,\lambda>0$, $q,r\geq 0$.
\end{enumerate}
\end{example}

\begin{remark}
The auxiliary function $\psi(s,t) = (s^p + r t^q)^\lambda$ generalizes several common forms:
\begin{itemize}
\item For $p = \lambda = 1$, $r = q = 0$: $\psi$ reduces to $d(Sx,Ty)$
\item For $r = 0$: $\psi$ becomes a power-scaling as in \cite{Rhoades1977}
\end{itemize}
This flexibility unifies prior work under one framework.
\end{remark}

\begin{theorem}[\cite{Choudhury2000}]
Let $T:X\to X$ be a selfmap of a complete metric space $X$. Assume that there exist $\psi\in\Psi$, $0<r<1$ and $0<s\leq 1$ such that

\[
\psi(d(Tx,Ty),d(x,Tx)) + \psi(d(y,Ty),d(y,T^{2}x)) \leq r\psi(d(x,y),d(x,Tx)) + s\psi(d(y,Ty),d(y,Tx)),
\]
where $x,y\in X$. Then $T$ has a unique fixed point in $X$.
\end{theorem}

Let 

\[
\Phi=\Bigg\{\phi\,|\,\phi:\mathbb{R}_{+}\to\mathbb{R}_{+} \text{ satisfying }
\begin{cases}
\text{(i) $\phi$ is monotonically increasing and continuous;} \\
\text{(ii) $0<\phi(t)<t$ for $t>0$}
\end{cases}
\Bigg\}
\]

\begin{definition}[\cite{Jungck1986}]
Two mappings $A$ and $T$ of a metric space $(X,d)$ are compatible if $\lim_{n\to\infty}d(ATx_{n},TAx_{n})=0$ whenever $\{x_{n}\}$ is a sequence in $X$ such that $\lim_{n\to\infty}Ax_{n}=\lim_{n\to\infty}Tx_{n}=t$ for some $t\in X$.
\end{definition}

In 1991, \cite{Antony1991} introduced the concept of weak compatibility of a pair selfmaps of a metric space as a generalization of compatibility, as follows:

\begin{definition}[\cite{Antony1991}]
Let $A$ and $T$ be selfmaps on a metric space $(X,d)$. The ordered pair $(T,A)$ is said to be weak compatible if $\lim_{n\to\infty}ATx_{n}=Tt$ whenever $\{x_{n}\}$ is a sequence in $X$ such that $\lim_{n\to\infty}Ax_{n}=\lim_{n\to\infty}Tx_{n}=t$ and $\lim_{n\to\infty}TAx_{n}=\lim_{n\to\infty}TTx_{n}=Tt$ for some $t\in X$.
\end{definition}

\begin{example}[\cite{Antony1991}] \label{ex:1.8}
The following example shows that two selfmaps of a metric space can be weak compatible without being compatible.

Let $X=\mathbb{R}$ with the usual metric. Define $A,T:X\to X$ by

\[
T(x)=\left\{\begin{array}{ll}
\frac{10}{32} & \text{if } x<\frac{3}{8} \\ 
\frac{3}{8} & \text{if } \frac{3}{8}\leq x<\frac{1}{2} \\
1 & \text{if } x\geq\frac{1}{2}
\end{array}\right.
\quad\text{and}\quad 
A(x)=\left\{\begin{array}{ll}
\frac{11}{32} & \text{if } x<\frac{3}{8} \\ 
\frac{1+x}{4} & \text{if } \frac{3}{8}\leq x<\frac{1}{2} \\ 
\frac{1+x}{2} & \text{if } x\geq\frac{1}{2}
\end{array}\right.
\]
We observe that $A$ and $T$ are not compatible while the ordered pair $(T,A)$ is weak compatible.
\end{example}

\begin{theorem}[\cite{Sastry2003}] \label{thm:1.9}
Let $(X,d)$ be a complete metric space and $A$ and $T$ be selfmaps on $X$ such that $(T,A)$ is weak compatible and satisfy: for some $\psi\in\Psi$, $\phi\in\Phi$ and for all $x,y\in X$,
\begin{align*}
    \psi(d(Tx,Ty),d(Tx,Ty)) \leq &\phi(\max\{\psi(d(Ax,Ay),d(Tx,Ax)), \psi(d(Ax,Ay),d(Ty,Ay)), \\
&\psi(d(Tx,Ax),d(Ty,Ay)), \psi(d(Ty,Ay),d(Tx,Ax)), \\
&\psi(d(Tx,Ay),d(Tx,Ax)), 
\psi(d(Tx,Ay),d(Ty,Ay))\}).
\end{align*}

If $T(X)\subset A(X)$ and $T$ is continuous, then $A$ and $T$ have a unique common fixed point.
\end{theorem}

\begin{definition}[\cite{Jungck1994}]
Two selfmaps $A$ and $T$ of a metric space $(X,d)$ are called weakly compatible if they commute at their coincidence point. That is, for $x\in X$, if $Ax=Tx$ then $ATx=TAx$.
\end{definition}

We observe that every compatible pair of selfmaps is weakly compatible but every weakly compatible pair of selfmaps need not be compatible \cite{Singh1997}.

We also observe that every weak compatible pair of selfmaps is weakly compatible ( Lemma 3.5D of \cite{Murthy1998}). But its converse need not be true as illustrated by Example \ref{ex:1.8}. In Example \ref{ex:1.8}, we observe that the ordered pair $(A,T)$ is not weak compatible while the pair of selfmaps $(A,T)$ is weakly compatible.

\begin{lemma}[\cite{KrishnaMurthy2004}] \label{lem:1.11}
Let $\psi:\mathbb{R}_{+}\to\mathbb{R}_{+}$ be upper semi-continuous and $\psi(t)<t$ for all $t\in\mathbb{R}_{+}-\{0\}$. Then there exists a monotone increasing continuous function $\phi:\mathbb{R}_{+}\to\mathbb{R}_{+}$ such that $\psi(t)\leq\phi(t)$ for all $t\in\mathbb{R}_{+}$ and $\phi(t)<t$ for all $t\in\mathbb{R}_{+}-\{0\}$.
\end{lemma}

%\section*{Auxiliary Functions of Two Variables}
The goal of this work is to prove a common fixed point theorem for two pairs of weakly compatible maps without requiring continuity, and under minimal completeness assumptions, using two-variable auxiliary functions. This general framework unifies and extends various known fixed point results including those of \cite{AlThagafi2006},  \cite{Jungck1994}, \cite{Murthy1998}, \cite{Pant1994}, \cite{Zhang2008}, and others.

\section{Main Results}

\begin{theorem}\label{thm:main}
Let $(X,d)$ be a metric space and $K$ be a nonempty subset of $X$. Let $A,B,S,T$ be selfmaps on $K$. Suppose there exist $\psi\in\Psi$, $\phi\in\Phi$ and such that for all $x,y\in K$:

\begin{align} \label{eq:2.1}
    \psi(d(Sx,Ty),d(Sx,Ty)) \leq \phi(&\max\{\psi(d(Ax,By),d(Ax,Sx)),\psi(d(Ax,By),d(By,Ty)),\nonumber \\
 &\psi(d(Ax,Sx),d(By,Ty)), \psi(d(By,Ty),d(Ax,Sx)),\\
 &\min\{\psi(d(By,Sx),d(Ax,Sx)),
\psi(d(Ax,Ty),d(By,Ty))\},  \nonumber\\
&\min\{\psi(d(By,Sx),d(By,Ty)), 
\psi(d(Ax,Ty),d(Ax,Sx))\}\}). \nonumber
\end{align}
Assume also that the pairs of selfmaps $(A,S)$ and $(B,T)$ are weakly compatible; $\overline{T(K)}\subset A(K)$ and $\overline{S(K)}\subset B(K)$. If either $\overline{T(K)}$ or $\overline{S(K)}$ or $A(K)$ or $B(K)$ is complete, then $A$, $B$, $S$ and $T$ have a unique common fixed point.
\end{theorem}

\begin{proof}
Let $x_{0}\in K$. Since $T(K)\subseteq\overline{T(K)}\subseteq A(K)$, we can find $x_{1}\in K$ such that $Tx_{0}=Ax_{1}$ and since $S(K)\subseteq\overline{S(K)}\subseteq B(K)$, we can find $x_{2}\in K$ such that $Sx_{1}=Bx_{2}$. Continuing this way,  We construct sequences  $\{x_{n}\}$ and $\{y_{n}\}$ in $K$ inductively as follows:
\[
y_{2n} = Tx_{2n} = Ax_{2n+1} \quad \text{and} \quad y_{2n+1} = Sx_{2n+1} = Bx_{2n+2}
\]
for $n=0,1,2,\cdots$.
Now we claim that the sequence $\{y_{n}\}_{n=1}^{\infty}$ is Cauchy in $K$.

\textbf{Case (i):} $y_{n}=y_{n+1}$ for some $n$.

Without any loss of generality, we assume that $n$ is even. Then $n=2m$ for some $m\in \mathbb{Z}_{+}$. So $y_{2m}=y_{2m+1}$.

Thus, we have
\begin{align*}
\psi(d(y_{n+1},y_{n+2}),d(y_{n+1},y_{n+2})) &= \psi(d(y_{2m+1},y_{2m+2}),d(y_{2m+1},y_{2m+2})) \\
&= \psi(d(Sx_{2m+1},Tx_{2m+2}),d(Sx_{2m+1},Tx_{2m+2})) \\
&\leq \phi(\max\{\psi(d(Ax_{2m+1},Bx_{2m+2}),d(Ax_{2m+1},Sx_{2m+1})), \\
&\quad \psi(d(Ax_{2m+1},Bx_{2m+2}),\psi(d(Bx_{2m+2}, Tx_{2m+2})),\\
& \psi(d(Ax_{2m+1}, Sx_{2m+1}), d(Bx_{2m+2}, Tx_{2m+2})), \\
&\psi(d(Bx_{2m+2}, Tx_{2m+2}), d(Ax_{2m+1}, Sx_{2m+1})), \\
&\min\{\psi(d(Bx_{2m+2}, Sx_{2m+1}), d(Ax_{2m+1}, Sx_{2m+1})), \\
&\quad \psi(d(Ax_{2m+1}, Tx_{2m+2}), d(Bx_{2m+2}, Tx_{2m+2}))\}, \\
&\min\{\psi(d(Bx_{2m+2}, Sx_{2m+1}), d(Bx_{2m+2}, Tx_{2m+2})), \\
&\quad \psi(d(Ax_{2m+1}, Tx_{2m+2}), d(Ax_{2m+1}, Sx_{2m+1}))\} \\
&= \phi(\max\{\psi(d(y_n, y_{n+1}), d(y_n, y_{n+1})), \\
&\quad \psi(d(y_n, y_{n+1}), d(y_{n+1}, y_{n+2})), \quad \psi(d(y_n, y_{n+1}), d(y_{n+1}, y_{n+2})),\\
&\quad \psi(d(y_{n+1}, y_{n+2}), d(y_n, y_{n+1})), \\
&\quad \min\{\psi(d(y_{n+1}, y_{n+1}), d(y_n, y_{n+1})), \quad \psi(d(y_n, y_{n+2}), d(y_{n+1}, y_{n+2}))\},\\
&\quad \min\{\psi(d(y_{n+1}, y_{n+1}), d(y_{n+1}, y_{n+2})), \quad \psi(d(y_n, y_{n+2}), d(y_n, y_{n+1}))\}\\
&\leq \phi(\max\{0, \psi(0, d(y_{n+1}, y_{n+2})), \psi(d(y_{n+1}, y_{n+2}), 0)\}) \\
&\leq \phi(\psi(d(y_{n+1}, y_{n+2}), d(y_{n+1}, y_{n+2}))).
\end{align*}

This implies that $\psi(d(y_{n+1}, y_{n+2}), d(y_{n+1}, y_{n+2})) = 0$ and hence $d(y_{n+1}, y_{n+2}) = 0$ (By Remark \ref{rem:1.3}).

Hence, $y_{n+1} = y_{n+2}$. Repeating this procedure inductively, we show that $y_n = y_{n+k}$ for $k \geq 1$. Hence, $\{y_m\}_{m\geq n}$ is a constant sequence and therefore is Cauchy in $K$.

\textbf{Case (ii):} $y_n \neq y_{n+1}$, $n = 1, 2, \dotsc$.

Let $\alpha_n = d(y_n, y_{n+1})$ for all $n = 0, 1, 2, \dotsc$.

Now consider
\begin{align*}
\psi(\alpha_{2n+1}, \alpha_{2n+1}) &= \psi(d(y_{2n+1}, y_{2n+2}), d(y_{2n+1}, y_{2n+2})) \\
&= \psi(d(Sx_{2n+1}, Tx_{2n+2}), d(Sx_{2n+1}, Tx_{2n+2})) \\
&\leq \phi(\max\{\psi(d(Ax_{2n+1}, Bx_{2n+2}), d(Ax_{2n+1}, Sx_{2n+1})), \psi(d(Ax_{2n+1}, Bx_{2n+2}), d(Bx_{2n+2}, Tx_{2n+2})),\\
&\quad \quad \psi(d(Ax_{2n+1}, Sx_{2n+1}), d(Bx_{2n+2}, Tx_{2n+2})), \psi(d(Bx_{2n+2}, Tx_{2n+2}), d(Ax_{2n+1}, Sx_{2n+1})),\\
&\quad \quad \min\{\psi(d(Bx_{2n+2}, Sx_{2n+1}), d(Ax_{2n+1}, Sx_{2n+1})), \psi(d(Ax_{2n+1}, Tx_{2n+2}), d(Bx_{2n+2}, Tx_{2n+2}))\}, \\
&\quad \quad \min\{\psi(d(Bx_{2n+2}, Sx_{2n+1}), d(Bx_{2n+2}, Tx_{2n+2})), \psi(d(Ax_{2n+1}, Tx_{2n+2}), d(Ax_{2n+1}, Sx_{2n+1}))\}\\
&= \phi(\max\{\psi(d(y_{2n}, y_{2n+1}), d(y_{2n}, y_{2n+1})), \psi(d(y_{2n}, y_{2n+1}), d(y_{2n+1}, y_{2n+2})),\\
&\quad \quad \psi(d(y_{2n}, y_{2n+1}), d(y_{2n+1}, y_{2n+2})), \psi(d(y_{2n+1}, y_{2n+2}), d(y_{2n}, y_{2n+1})), \\
&\quad \quad \min\{\psi(d(y_{2n+1}, y_{2n+1}), d(y_{2n}, y_{2n+1})), \psi(d(y_{2n}, y_{2n+2}), d(y_{2n+1}, y_{2n+2}))\}, \\
&\quad \quad \min\{\psi(d(y_{2n+1}, y_{2n+1}), d(y_{2n+1}, y_{2n+2})), \psi(d(y_{2n}, y_{2n+2}), d(y_{2n}, y_{2n+1}))\} \\
&\leq \phi(\max\{\psi(\alpha_{2n}, \alpha_{2n}), \psi(\alpha_{2n}, \alpha_{2n+1}), \psi(\alpha_{2n+1}, \alpha_{2n}),\psi(0, \alpha_{2n}), \psi(0, \alpha_{2n+1})\}).
\end{align*}

Now if $\alpha_{2n} < \alpha_{2n+1}$, then $\psi(\alpha_{2n+1}, \alpha_{2n+1}) \leq \phi(\psi(\alpha_{2n+1}, \alpha_{2n+1}))$, which is a contradiction. Hence, $\alpha_{2n+1} \leq \alpha_{2n}$ and
\begin{equation}\label{eq:2.1.2}
\psi(\alpha_{2n+1}, \alpha_{2n+1}) \leq \phi(\psi(\alpha_{2n}, \alpha_{2n})), \quad n = 0, 1, 2, \dotsc.
\end{equation}

Similarly, we can show that
\begin{equation}\label{eq:2.1.3}
\psi(\alpha_{2n}, \alpha_{2n}) \leq \phi(\psi(\alpha_{2n-1}, \alpha_{2n-1})), \quad n = 1, 2, 3, \dotsc.
\end{equation}

Hence, from \eqref{eq:2.1.2} and \eqref{eq:2.1.3}, we get
\begin{equation}\label{eq:2.1.4}
\psi(\alpha_{n+1}, \alpha_{n+1}) \leq \phi(\psi(\alpha_{n}, \alpha_{n})), \quad n = 0, 1, 2, \dotsc.
\end{equation}

Since $\psi$ is monotone increasing in each of its variables, we have
\begin{equation}\label{eq:2.1.6}
\alpha_{2n+1} \leq \alpha_{2n}, \quad n = 0, 1, 2, \dotsc.
\end{equation}

Thus, the sequence $\{\alpha_n\}$ decreases, say, to $\alpha (\geq 0)$. Since $\phi$ and $\psi$ are continuous, letting $n \to \infty$, from \eqref{eq:2.1.4}, we get
\[
\psi(\alpha, \alpha) \leq \phi(\psi(\alpha, \alpha)),
\]
so that $\psi(\alpha, \alpha) = 0$ and hence $\alpha = 0$. Therefore,
\begin{equation}\label{eq:2.1.7}
\lim_{n\to\infty} d(y_n, y_{n+1}) = 0.
\end{equation}

We now claim that $\{y_n\}$ is Cauchy. By \eqref{eq:2.1.6} and \eqref{eq:2.1.7}, it is sufficient to show that $\{y_{2n}\}$ is Cauchy. Otherwise, there exists an $\epsilon > 0$ and there exist sequences $\{m_k\}$ and $\{n_k\}$ with $m_k > n_k > k$ such that
\begin{equation}\label{eq:2.1.8}
d(y_{2m_k}, y_{2n_k}) \geq \epsilon \quad \text{and} \quad d(y_{2m_k-2}, y_{2n_k}) < \epsilon.
\end{equation}

Hence,
\begin{equation}\label{eq:2.1.9}
\epsilon \leq \liminf_{k\to\infty} d(y_{2m_k}, y_{2n_k}).
\end{equation}

Now for each positive integer $k$, by triangle inequality we get,
\begin{equation*}
d(y_{2m_k}, y_{2n_k}) \leq d(y_{2m_k}, y_{2m_k-1}) + d(y_{2m_k-1}, y_{2m_k-2}) + d(y_{2m_k-2}, y_{2n_k}).
\end{equation*}

On taking limit supremum of both sides, as $k \to \infty$, we get
\begin{equation}\label{eq:2.1.10}
\limsup_{k\to\infty} d(y_{2m_k}, y_{2n_k}) \leq \epsilon.
\end{equation}

Hence from \eqref{eq:2.1.9} and \eqref{eq:2.1.10}, we have
\begin{equation}\label{eq:2.1.11}
\lim_{k\to\infty} d(y_{2m_k}, y_{2n_k}) = \epsilon.
\end{equation}

Now
\begin{align*}
d(y_{2m_k}, y_{2n_k-1}) &\leq d(y_{2m_k}, y_{2n_k}) + d(y_{2n_k}, y_{2n_k-1});
\end{align*}
On taking limit supremum, as $k \to \infty$, we get
\begin{equation}\label{eq:2.1.12}
\limsup_{k\to\infty} d(y_{2m_k}, y_{2n_k-1}) \leq \epsilon.
\end{equation}

Again we have
\begin{align*}
d(y_{2m_k}, y_{2n_k}) &\leq d(y_{2m_k}, y_{2n_k-1}) + d(y_{2n_k-1}, y_{2n_k});
\end{align*}
On taking limit infimum, as $k \to \infty$, we get
\begin{equation}\label{eq:2.1.13}
\epsilon \leq \liminf_{k\to\infty} d(y_{2m_k}, y_{2n_k-1}).
\end{equation}

From \eqref{eq:2.1.12} and \eqref{eq:2.1.13}, we have
\begin{equation}\label{eq:2.1.14}
\lim_{k\to\infty} d(y_{2m_k}, y_{2n_k-1}) = \epsilon.
\end{equation}

Similarly, we can show that
\begin{equation}\label{eq:2.1.15}
\lim_{k\to\infty} d(y_{2m_k+1}, y_{2n_k}) = \epsilon
\end{equation}
and
\begin{equation}\label{eq:2.1.16}
\lim_{k\to\infty} d(y_{2m_k+1}, y_{2n_k-1}) = \epsilon.
\end{equation}

Now consider
\begin{align*}
\psi(d(y_{2m_k+1}, y_{2n_k}), d(y_{2m_k+1}, y_{2n_k})) &= \psi(d(Sx_{2m_k+1}, Tx_{2n_k}), d(Sx_{2m_k+1}, Tx_{2n_k})) \\
&\leq \phi(\max\{\psi(d(Ax_{2m_k+1}, Bx_{2n_k}), d(Ax_{2m_k+1}, Sx_{2m_k+1})), \\
&\quad \psi(d(Ax_{2m_k+1}, Bx_{2n_k}), d(Bx_{2n_k}, Tx_{2n_k})), \\
&\quad \psi(d(Ax_{2m_k+1}, Sx_{2m_k+1}), d(Bx_{2n_k}, Tx_{2n_k})), \\
&\quad \psi(d(Bx_{2n_k}, Tx_{2n_k}), d(Ax_{2m_k+1}, Sx_{2m_k+1})), \\
&\quad \min\{\psi(d(Bx_{2n_k}, Sx_{2m_k+1}), d(Ax_{2m_k+1}, Sx_{2m_k+1})), \\
&\quad \quad \psi(d(Ax_{2m_k+1}, Tx_{2n_k}), d(Bx_{2n_k}, Tx_{2n_k}))\}, \\
&\quad \min\{\psi(d(Bx_{2n_k}, Sx_{2m_k+1}), d(Bx_{2n_k}, Tx_{2n_k})), \\
&\quad \quad \psi(d(Ax_{2m_k+1}, Tx_{2n_k}), d(Ax_{2m_k+1}, Sx_{2m_k+1}))\} \\
&= \phi(\max\{\psi(d(y_{2m_k}, y_{2n_k-1}), d(y_{2m_k}, y_{2m_k+1})), \psi(d(y_{2m_k}, y_{2n_k-1}), d(y_{2n_k-1}, y_{2n_k})),\\
&\quad \min\{\psi(d(y_{2n_k-1},y_{2m_k+1}),d(y_{2m_k},y_{2m_k+1})), \psi(d(y_{2m_k},y_{2n_k}),d(y_{2n_k-1},y_{2n_k}))\},\\
&\quad \min\{\psi(d(y_{2n_k-1},y_{2m_k+1}),d(y_{2n_k-1},y_{2n_k})),\psi(d(y_{2m_k},y_{2n_k}),d(y_{2m_k},y_{2m_k+1}))\}
\end{align*}

Letting $n\to\infty$, by the continuity of $\phi$ and $\psi$, and from \eqref{eq:2.1.7}, \eqref{eq:2.1.11}, \eqref{eq:2.1.14}, \eqref{eq:2.1.15} and \eqref{eq:2.1.16}, we get

\begin{align*}
\psi(\varepsilon,\varepsilon) &\leq \phi(\max\{\psi(\varepsilon,0),\psi(\varepsilon,0),\psi(0,0), \psi(0,0), \\
&\quad \min\{\psi(\varepsilon,0),\psi(\varepsilon,0)\}, \min\{\psi(\varepsilon,0),\psi(\varepsilon,0)\}\}) \\
&= \phi(\psi(\varepsilon,0)) \\
&\leq \phi(\psi(\varepsilon,\varepsilon)), \text{ a contradiction}.
\end{align*}

Hence, $\{y_{2n}\}$ is a Cauchy sequence in $K$ and hence $\{y_n\}$ is Cauchy in $K$.

Hence, the subsequences $\{Tx_{2n}\}$, $\{Sx_{2n+1}\}$, $\{Ax_{2n}\}$ and $\{Bx_{2n+2}\}$ are Cauchy in $K$.

Assume that $\overline{T(K)}$ is complete. Then there exists a $z\in\overline{T(K)}$ such that
\begin{equation}\label{eq:2.1.17}
\lim_{n\to\infty} Tx_{2n} = z
\end{equation}
and hence,
\begin{equation}\label{eq:2.1.18}
\lim_{n\to\infty} Ax_{2n+1} = \lim_{n\to\infty} Bx_{2n+2} = \lim_{n\to\infty} Sx_{2n+1} = z
\end{equation}
and $z\in\overline{S(K)}$.

Since $\overline{T(K)} \subset A(K)$ and $\overline{S(K)} \subset B(K)$, there exist $u,v\in K$ such that
\begin{equation}\label{eq:2.1.19}
z = Au \text{ and } z = Bv.
\end{equation}

Now taking $x = u$ and $y = x_{2n}$ in \eqref{eq:2.1}, we have
\begin{align*}
\psi(d(Su,Tx_{2n}),d(Su,Tx_{2n})) &\leq \phi(\max\{\psi(d(Au,Bx_{2n}),d(Au,Su)),\psi(d(Au,Bx_{2n}), d(Bx_{2n},Tx_{2n})),  \\
&\quad \psi(d(Au,Su),d(Bx_{2n},Tx_{2n})), \psi(d(Bx_{2n},Tx_{2n}),d(Au,Su)),\\
&\quad \min\{\psi(d(Bx_{2n},Su),d(Au,Su)), \psi(d(Au,Tx_{2n}),d(Bx_{2n},Tx_{2n}))\},\\
&\quad \min\{\psi(d(Bx_{2n},Su), d(Bx_{2n},Tx_{2n})), \psi(d(Au,Tx_{2n}),d(Au,Su))\})
\end{align*}

Letting $n\to\infty$, by the continuity and monotone increasing property of $\phi$ and $\psi$ and using \eqref{eq:2.1.17}, \eqref{eq:2.1.18} and \eqref{eq:2.1.19}, we get
\[
\psi(d(Su,z),d(Su,z)) \leq \phi(\psi(d(Su,z),d(Su,z))).
\]
This implies that $\psi(d(Su,z),d(Su,z)) = 0$ and hence by Remark \ref{rem:1.3} we have $d(z,Su) = 0$.

Hence, $z = Su$.

Similarly in \eqref{eq:2.1}, by taking $x = x_{2n+1}$ and $y = v$ and letting $n\to\infty$, we get $z = Tv$.

Hence, $Su = Au = z = Tv = Bv$.

Since the pairs of mappings $(A,S)$ and $(B,T)$ are weakly compatible mappings, we have $ASu = SAu$ and $BTv = TBv$ and hence
\[
Az = Sz \text{ and } Bz = Tz.
\]

Hence, $z$ is the point of coincidence of $A$ and $S$, and $B$ and $T$.

We now claim that $Az = Sz = z$. Consider
\begin{align*}
\psi(d(Sz,z),d(Sz,z)) &= \psi(d(Sz,Tu),d(Sz,Tu)) \\
&\leq \phi(\max\{\psi(d(Az,Bu),d(Az,Sz)), \psi(d(Az,Bu),d(Bu,Tu)), \\
&\quad \psi(d(Az,Sz),d(Bu,Tu)), \psi(d(Bu,Tu),d(Az,Sz)), \\
&\quad \min\{\psi(d(Bu,Sz),d(Az,Sz)), \psi(d(Az,Tu),d(Bu,Tu))\}, \\
&\quad \min\{\psi(d(Bu,Sz),d(Bu,Tu)), \psi(d(Az,Tu),d(Az,Sz))\}\}) \\
&= \phi(\psi(d(Sz,z),0)) \\
&\leq \phi(\psi(d(Sz,z),d(Sz,z))).
\end{align*}

This implies that $\psi(d(Sz,z),d(Sz,z)) = 0$ and hence $d(Sz,z) = 0$. Therefore,
\[
Sz = z.
\]
Hence, $Az = Sz = z$.

Similarly, by taking $x = u$ and $y = z$ in \eqref{eq:2.1} we get $Tz = z$. Hence,
\[
Bz = Tz = z.
\]

Thus, $z$ is a common fixed point of $A$, $B$, $S$ and $T$. The uniqueness of $z$ follows from the inequality \eqref{eq:2.1}.

Similarly, by symmetry we prove for the case if $\overline{S(K)}$ is complete.

Now assume $A(K)$ is complete. As $\overline{T(K)} \subset A(K)$ and $\overline{T(K)}$ is closed, $\overline{T(K)}$ is complete. Hence by the above argument, there exists a point $z \in \overline{T(K)}$ such that
\[
Az = Bz = Sz = Tz = z.
\]

Similarly, if we assume $B(K)$ is complete, the proof follows by symmetry.

Hence the theorem.
\end{proof}

The following corollaries are immediate from Theorem \ref{thm:main}.

\begin{corollary}\label{cor:2.2}
Let $(X,d)$ be a metric space and $K$ be a nonempty subset of $X$. Let $A,B,S,T$ be selfmaps on $K$. Assume that there exist $r\in[0,1)$ and $\psi\in\Psi$, and for all $x,y\in X$,
\begin{align*}
    \psi(d(Sx,Ty),d(Sx,Ty)) \leq r &\max\{\psi(d(Ax,By),d(Ax,Sx)), \psi(d(Ax,By),d(By,Ty)),\\
 &\psi(d(Ax,Sx),d(By,Ty)), \psi(d(By,Ty),d(Ax,Sx)),\\
 &\min\{\psi(d(By,Sx),d(Ax,Sx)),\psi(d(Ax,Ty),d(By,Ty))\}, \\
 &\min\{\psi(d(By,Sx),d(By,Ty)), \psi(d(Ax,Ty),d(Ax,Sx))\}\}.
\end{align*}
Assume also that the pairs of selfmaps $(A,S)$ and $(B,T)$ are weakly compatible; $T(K)\subset A(K)$ and $\overline{S(K)}\subset B(K)$. If either $T(K)$ or $\overline{S(K)}$ or $A(K)$ or $B(K)$ is complete, then $A$, $B$, $S$ and $T$ have a unique common fixed point.
\end{corollary}

\begin{proof}
Follows from Theorem \ref{thm:main} by choosing $\phi(t)=rt$, for some $r\in[0,1)$ and for all $t\in \mathbb{R}_{+}$.
\end{proof}

\begin{corollary}\label{cor:2.3}
Let $(X,d)$ be a metric space and $K$ be a nonempty subset of $X$. Let $A$ and $T$ be selfmaps on $K$. Assume that there exist $\psi\in\Psi$ and $\phi\in\Phi$, and for all $x,y\in X$,

\begin{align*}
    \psi(d(Tx,Ty),d(Tx,Ty)) \leq &\phi(\max\{\psi(d(Ax,Ay),d(Ax,Tx)), \psi(d(Ax,Ay),d(Ty,Ty)),\\
 &\psi(d(Ax,Tx),d(Ay,Ty)), \psi(d(Ay,Ty),d(Ax,Tx)),\\
 &\min\{\psi(d(Ay,Tx),d(Ax,Tx)), \psi(d(Ax,Ty),d(Ay,Ty))\},\\
 &\min\{\psi(d(Ay,Tx),d(Ay,Ty)), \psi(d(Ax,Ty),d(Ax,Tx))\}\}).
\end{align*}
Assume also that the pair of selfmaps $(A,T)$ is weakly compatible; $\overline{T(K)}\subset A(K)$. If either $T(K)$ or $A(K)$ is complete, then $A$ and $T$ have a unique common fixed point.
\end{corollary}

\begin{proof}
Follows from Theorem \ref{thm:main} by choosing $S=T$ and $B=A$.
\end{proof}

\begin{remark}\label{rem:2.4}
Corollary \ref{cor:2.3} improves Theorem \ref{thm:1.9}, for, in Theorem \ref{thm:1.9} $X$ is complete metric space, the map $T$ is continuous and the ordered pair $(T,A)$ is weak compatible. Here we have replaced the completeness of $X$ by completeness of either one of $A(X)$, $B(X)$, $\overline{S(X)}$ or $\overline{T(X)}$. Also in Corollary \ref{cor:2.3}, no continuity assumption of the mappings is used. Further, the weak compatibility of the ordered pair of mappings $(T,A)$ is relaxed by weakly compatibility of $A$ and $T$.
\end{remark}

\begin{theorem}\label{thm:2.5}
Let $(X, d)$ be a metric space and $K$ be a nonempty subset of $X$. Let $A, B, S, T$ be selfmaps on $K$. Assume that there exists $\phi \in \Phi$ such that
\begin{align*}
    [(d(Sx, Ty))^p +& r (d(Sx, Ty))^q]^\lambda \leq \phi(\max\{[(d(Ax, By))^p + r (d(Ax, Sx))^q]^\lambda, \\
&[(d(Ax, By))^p + r (d(By, Ty))^q]^\lambda, [(d(Ax, Sx))^p + r (d(By, Ty))^q]^\lambda \\
&[(d(By, Ty))^p + r (d(Ax, Sx))^q]^\lambda, \\
&\min\{[(d(By, Sx))^p + r (d(Ax, Sx))^q]^\lambda, [(d(Ax, Ty))^p + r(d(By, Ty))^q]^\lambda\} \\
&\min\{[(d(By, Sx))^p + r (d(By, Ty))^q]^\lambda, [(d(Ax, Ty))^p + r (d(Ax, Sx))^q]^\lambda\}\})
\end{align*}
for all $x, y \in K$, for some $p, \lambda > 0$ and $r, q \geq 0$.

Assume also that the pairs of selfmaps $(A, S)$ and $(B, T)$ are weakly compatible; $T(K) \subset A(K)$ and $S(K) \subset B(K)$. If either $T(K)$ or $S(K)$ or $A(K)$ or $B(K)$ is complete, then $A$, $B$, $S$ and $T$ have a unique common fixed point.
\end{theorem}

\begin{proof}
Define $\psi(a, b) = (a^p + r b^q)^\lambda$ in Theorem \ref{thm:main}.
\end{proof}

\begin{remark}\label{rem:2.6}
\begin{enumerate}
\item[(i)] By Lemma \ref{lem:1.11}, for $S = T$, $p = \lambda = 1$ and $r = q = 0$, Theorem 2.1 of \cite{Zhang2008} follows as a corollary to Theorem \ref{thm:2.5} under the additional assumption that $\phi : \mathbb{R}_{+} \to \mathbb{R}_{+}$ is upper semi continuous from the right.

\item[(ii)] For $\phi(t) = rt$, for some $r \in [0, 1)$, $p = \lambda = 1$ and $r = q = 0$, Theorem \ref{thm:2.5} modifies Corollary 3.2 of \cite{Jungck1986} and Corollary 2.1 of \cite{Jungck1994} if we take $X = K$.

\item[(iii)] For $\phi(t) = rt$, for some $r \in [0, 1)$, $p = \lambda = 1$ and $r = q = 0$, Theorem \ref{thm:2.5} extends Theorem 2.1 of \cite{AlThagafi2006}.

\item[(iv)] For $p = \lambda = 1$ and $r = q = 0$, $B = A$ and $S = T$, Theorem 1 of \cite{Pant1994} and Theorem 3.6 of \cite{Murthy1998} are corollaries to Theorem \ref{thm:2.5} under the additional assumption that $\phi$ is monotone increasing in their results.

\item[(v)] If we choose $S = T$ and $B = A = I$, the identity mapping, in Theorem \ref{thm:2.5}, for various choices of $p, q, r$ and $\lambda$ in Theorem \ref{thm:2.5} we can obtain the results of various authors as corollaries including Banach contraction theorem. For example, by taking $S = T$, $B = A = I$, the identity mapping, $p = 1$, $r = 0$, $q = 1$, $\lambda = 1$ in Theorem \ref{thm:2.5}, we obtain a generalization of some results of \cite{Bianchini1972}, \cite{Kannan1968}, \cite{Reich1971} and \cite{Rhoades1977}.
\end{enumerate}
\end{remark}

\section{Conclusion}
In this work, we established novel common fixed point theorems for weakly compatible mappings using auxiliary functions of two variables. Our results significantly generalize existing theorems by:  

\begin{itemize}  
\item Relaxing the continuity and completeness assumptions, requiring only the completeness of \emph{one} image subset (e.g., $\overline{T(K)}$ or $\overline{S(K)}$).  
\item Employing a flexible auxiliary function $\psi \in \Psi$ to unify diverse contraction conditions, as illustrated in Example 1.4 and Theorem 2.5.  
\item Demonstrating that weak compatibility (Definition 1.10) suffices, replacing stricter notions like compatibility or weak compatibility (Definition 1.7).  
\end{itemize}  

Theorems 2.1 and 2.5, along with their corollaries, extend and refine prior work from \cite{AlThagafi2006, Jungck1986, Jungck1994, Murthy1998, Pant1994, Zhang2008} while maintaining minimal assumptions. 

\begin{thebibliography}{10}

\bibitem{AlThagafi2006}
M.~A. Al-Thagafi and N.~Shahzad.
\newblock Noncommuting selfmaps and invariant approximations.
\newblock {\em Nonlinear Analysis}, 64:2778--2786, 2006.

\bibitem{Babu2004}
G.~V.~R. Babu.
\newblock Generalization of fixed point theorems related to the diameter of orbits by using a control function.
\newblock {\em Tamkang Journal of Mathematics}, 35(2):159--168, 2004.

\bibitem{Babu2001}
G.~V.~R. Babu and Shaik Ismail.
\newblock A fixed point theorem by altering distances.
\newblock {\em Bulletin of the Calcutta Mathematical Society}, 93(5):393--398, 2001.

\bibitem{Babu2007}
G.~V.~R. Babu, B.~Lalitha, and M.~L. Sandhya.
\newblock Common fixed point theorems involving two generalized altering distance functions in four variables.
\newblock {\em Proceedings of the Jangjeon Mathematical Society}, 10:83--93, 2007.

\bibitem{Babu2007b}
G.~V.~R. Babu and K.~N. V.~V. Vara~Prasad.
\newblock Common fixed point theorems for a pair of compatible maps using a control function.
\newblock {\em Proceedings of the Jangjeon Mathematical Society}, 10(1):55--70, 2007.

\bibitem{Bianchini1972}
R.~M.~T. Bianchini.
\newblock Su un problema di s. reich riguardante la teoria dei punti fissi.
\newblock {\em Bollettino dell'Unione Matematica Italiana}, 5:103, 1972.

\bibitem{Choudhury2005}
Binayak~S. Choudhury.
\newblock A common unique fixed point result in metric spaces involving generalized altering distances.
\newblock {\em Mathematical Communications}, 10:105--110, 2005.

\bibitem{Choudhury2000}
Binayak~S. Choudhury and P.~N. Dutta.
\newblock A unified fixed point result in metric spaces involving a two-variable function.
\newblock {\em FILOMAT}, 14:43--48, 2000.

\bibitem{Antony1991}
Antony Jessy.
\newblock {\em Studies in fixed points and quasi-gauges}.
\newblock PhD thesis, Indian Institute of Technology, Madras, India, 1991.

\bibitem{Jungck1986}
G.~Jungck.
\newblock Compatible mappings and common fixed points.
\newblock {\em International Journal of Mathematics and Mathematical Sciences}, 9(4):771--779, 1986.

\bibitem{Jungck1994}
G.~Jungck.
\newblock Compatible mappings and common fixed points (revised).
\newblock {\em International Journal of Mathematics and Mathematical Sciences}, 17(1):37--40, 1994.

\bibitem{Kannan1968}
R.~Kannan.
\newblock Some results on fixed points.
\newblock {\em Bulletin of the Calcutta Mathematical Society}, 60:71--76, 1968.

\bibitem{Khan1984}
M.~S. Khan, M.~Swaleh, and S.~Sessa.
\newblock Fixed point theorems by altering distances between the points.
\newblock {\em Bulletin of the Australian Mathematical Society}, 30:1--9, 1984.

\bibitem{KrishnaMurthy2004}
I.~S.~R. Krishna~Murthy.
\newblock {\em A study on the existence of unique common fixed points for self-maps on complete and compact metric spaces}.
\newblock PhD thesis, Andhra University, Visakhapatnam, India, 2004.

\bibitem{Murthy1998}
A.~S.~R. Murthy.
\newblock Some fixed point theorems in metric and quasi-gauge spaces, 1998.

\bibitem{Naidu2001}
S.~V.~R. Naidu.
\newblock Fixed point theorems by altering distances.
\newblock {\em Advances in Mathematical Sciences and Applications}, 11(1):1--16, 2001.

\bibitem{Naidu2003}
S.~V.~R. Naidu.
\newblock Fixed point theorems in metric spaces by altering distances.
\newblock {\em Czechoslovak Mathematical Journal}, 53(128):205--212, 2003.

\bibitem{Pant1994}
R.~P. Pant.
\newblock Common fixed points of noncommuting mappings.
\newblock {\em Journal of Mathematical Analysis and Applications}, 188:436--440, 1994.

\bibitem{Park1980}
S.~Park.
\newblock A unified approach to fixed points of contractive maps.
\newblock {\em Journal of the Korean Mathematical Society}, 16:95--105, 1980.

\bibitem{Pathak1994}
H.~K. Pathak and Rekha Sharma.
\newblock A note on fixed point theorems of khan, swaleh, and sessa.
\newblock {\em Mathematics Education}, 28(3):151--157, 1994.

\bibitem{Reich1971}
S.~Reich.
\newblock Some remarks concerning contraction mappings.
\newblock {\em Canadian Mathematical Bulletin}, 14:121, 1971.

\bibitem{Rhoades1977}
B.~E. Rhoades.
\newblock A comparison of various definitions of contractive mappings.
\newblock {\em Transactions of the American Mathematical Society}, 226:257--290, 1977.

\bibitem{Sastry1998}
K.~P.~R. Sastry and G.~V.~R. Babu.
\newblock Fixed point theorems in metric spaces by altering distances.
\newblock {\em Bulletin of the Calcutta Mathematical Society}, 90:175--182, 1998.

\bibitem{Sastry1999}
K.~P.~R. Sastry and G.~V.~R. Babu.
\newblock Some fixed point theorems by altering distances between the points.
\newblock {\em Indian Journal of Pure and Applied Mathematics}, 30(6):641--647, 1999.

\bibitem{Sastry2003}
K.~P.~R. Sastry and I.~S.~R. Krishna~Murthy.
\newblock Fixed point theorems through auxiliary functions.
\newblock {\em Applied Science Periodicals}, 5(1):52--57, 2003.

\bibitem{Sastry2000}
K.~P.~R. Sastry, S.~V.~R. Naidu, G.~V.~R. Babu, and G.~A. Naidu.
\newblock Generalization of common fixed point theorems by altering distances.
\newblock {\em Tamkang Journal of Mathematics}, 31:243--250, 2000.

\bibitem{Singh1997}
S.~L. Singh and S.~N. Mishra.
\newblock Remarks on jachymski's fixed point theorems for compatible maps.
\newblock {\em Indian Journal of Pure and Applied Mathematics}, 28:611--615, 1997.

\bibitem{Zhang2008}
S.~Zhang, Q.~Zhu, and Y.~Song.
\newblock Altering picard iterates for hybrid boyd-wong contractions.
\newblock {\em International Journal of Mathematical Analysis}, 2(12):563--568, 2008.

\end{thebibliography}


\end{document}
